\begin{table}[H]
\centering
\caption[\textbf{Morse Potentials}]{\textbf{Morse potentials.} {\normalfont
Morse potential representation given by $V(r) = D_e\left(1 - e^{-a(r-r_e)}\right)^2$,
used to describe bond interactions. Atom IDs correspond to those shown in Fig.~S2.
Here, SEC denotes selenium-containing systems, SUL denotes sulfur-containing systems,
and PRX refers to hydrogen peroxide ($\mathrm{H_2O_2}$).}}
\begin{threeparttable}
\begin{tabularx}{\textwidth}{lXXX}
\toprule
\textbf{Bond type} &
\makecell{\textbf{H--O (peroxide,} \\ \textbf{water) [$D_e$, $\alpha$, $r_e$]}} &
\makecell{\textbf{H--Se/H--S} \\ \textbf{[$D_e$, $\alpha$, $r_e$]}} &
\makecell{\textbf{HE1--OE1} \\ \textbf{[$D_e$, $\alpha$, $r_e$]}} \\
\midrule
SEC & 102, 2.26, 0.96 & 81, 1.87, 1.40\tnote{*} & 102, 2.26, 0.96 \\
SUL & 102, 2.26, 0.96 & 81, 1.87, 1.32\tnote{*} & 102, 2.26, 0.96 \\
\bottomrule
\end{tabularx}
\begin{tablenotes}
\item[*] The H--Se/H--S value differs between SEC and SUL; both values are intentionally retained for comparison.
\end{tablenotes}
\end{threeparttable}
\end{table}

%% TABLE 2 — Flat-bottom potential bound distances
\begin{table}[H]
\centering
\small
\caption[Flat-bottom potential bound distances]{\textbf{Flat-bottom potential bound distances.} {\normalfont Flat-bottom potential bound distances for reacting atom pairs, expressed as a constant $\pm 0.20$~\si{\angstrom} offset from the equilibrium distance.}}
\label{tab:flatbottom}
\begin{tabular}{lcc}
\toprule
\textbf{Pair} & \textbf{Equilibrium Distance (\si{\angstrom})} & \textbf{Bound Relative to Equilibrium (\si{\angstrom})} \\
\midrule
S--O  & 1.7  & $\pm 0.20$ \\
S--H  & 1.34 & $\pm 0.20$ \\
Se--O & 1.5  & $\pm 0.20$ \\
Se--H & 1.47 & $\pm 0.20$ \\
O--O  & 1.48 & $\pm 0.20$ \\
O--H  & 0.96 & $\pm 0.20$ \\
\bottomrule
\end{tabular}
\end{table}

%% TABLE 3 — Partial charges for the Cys system
\begin{table}[htbp]
\centering
\small
\caption[Partial charges for the Cys system]{\textbf{Partial charges for the Cys system.} {\normalfont Partial atomic charges in reactant state (RS) and product state (PS) for atoms in CYS, GLN, and hydrogen peroxide (\ce{H2O2}).}}
\label{tab:partial_charges}
\begin{tabular}{llcc}
\toprule
\textbf{Moiety} & \textbf{Atom} & \textbf{RS Charge} & \textbf{PS Charge} \\
\midrule
CYS   & SG   & $-0.335$ & $-0.9$    \\
CYS   & HG1  & $+0.155$ & $+0.418$  \\
CYS   & CB   & $+0.060$ & $-0.3$    \\
CYS   & HB1  & $+0.060$ & $+0.1$    \\
CYS   & HB2  & $+0.060$ & $+0.1$    \\
CYS   & CA   & $+0.140$ & $+0.14$   \\
CYS   & HA   & $+0.060$ & $+0.06$   \\
CYS   & N    & $-0.500$ & $-0.5$    \\
CYS   & HN   & $+0.300$ & $+0.3$    \\
\addlinespace
\ce{H2O2} & O1 & $-0.418$ & $-0.418$ \\
\ce{H2O2} & O2 & $-0.418$ & $-0.418$ \\
\ce{H2O2} & H1 & $+0.418$ & $+0.418$ \\
\ce{H2O2} & H2 & $+0.418$ & $+0.418$ \\
\addlinespace
GLN   & NE2  & $-0.760$ & $-0.4965$ \\
GLN   & HE21 & $+0.380$ & $+0.48$   \\
GLN   & HE22 & $+0.380$ & $+0.48$   \\
GLN   & CD   & $+0.500$ & $+0.2858$ \\
GLN   & OE1  & $-0.500$ & $-0.503$  \\
GLN   & CB   & $-0.120$ & $-0.12$   \\
GLN   & HB1  & $+0.060$ & $+0.06$   \\
GLN   & HB2  & $+0.060$ & $+0.06$   \\
GLN   & CG   & $-0.120$ & $+0.2157$ \\
GLN   & HG1  & $+0.060$ & $+0.06$   \\
GLN   & HG2  & $+0.060$ & $+0.06$   \\
\bottomrule
\end{tabular}
\end{table}

%% TABLE 4 — Partial charges for the Sec system
\begin{table}[htbp]
\centering
\small
\caption[Partial charges for the Sec system]{\textbf{Partial charges for the Sec system.} {\normalfont Partial atomic charges in reactant state (RS) and product state (PS) for atoms in SEC, GLN, and hydrogen peroxide (\ce{H2O2}).}}
\label{tab:partial_charges_sec}
\begin{tabular}{llcc}
\toprule
\textbf{Moiety} & \textbf{Atom} & \textbf{RS Charge} & \textbf{PS Charge} \\
\midrule
SEC   & SE   & $-0.335$ & $-0.9$    \\
SEC   & HG1  & $+0.155$ & $+0.418$  \\
SEC   & CB   & $+0.060$ & $-0.3$    \\
SEC   & HB1  & $+0.060$ & $+0.1$    \\
SEC   & HB2  & $+0.060$ & $+0.1$    \\
SEC   & CA   & $+0.140$ & $+0.14$   \\
SEC   & HA   & $+0.060$ & $+0.06$   \\
SEC   & N    & $-0.500$ & $-0.5$    \\
SEC   & HN   & $+0.300$ & $+0.3$    \\
\addlinespace
\ce{H2O2} & O1 & $-0.418$ & $-0.418$ \\
\ce{H2O2} & O2 & $-0.418$ & $-0.418$ \\
\ce{H2O2} & H1 & $+0.418$ & $+0.418$ \\
\ce{H2O2} & H2 & $+0.418$ & $+0.418$ \\
\addlinespace
GLN   & NE2  & $-0.760$ & $-0.4965$ \\
GLN   & HE21 & $+0.380$ & $+0.48$   \\
GLN   & HE22 & $+0.380$ & $+0.48$   \\
GLN   & CD   & $+0.500$ & $+0.2858$ \\
GLN   & OE1  & $-0.500$ & $-0.503$  \\
GLN   & CB   & $-0.120$ & $-0.12$   \\
GLN   & HB1  & $+0.060$ & $+0.06$   \\
GLN   & HB2  & $+0.060$ & $+0.06$   \\
GLN   & CG   & $-0.120$ & $+0.2157$ \\
GLN   & HG1  & $+0.060$ & $+0.06$   \\
GLN   & HG2  & $+0.060$ & $+0.06$   \\
\bottomrule
\end{tabular}
\end{table}

%% TABLE 5 — EVB Parameters
\begin{table}[htbp]
  \centering
  \small
  \caption[EVB Parameters]{\textbf{EVB parameters.} {\normalfont Mapped EVB parameters used for the WT reference calculations.}}
  \label{tab:evb_parameters}
  \begin{tabular}{lc}
    \toprule
    \textbf{Parameter} & \textbf{Value (kcal~mol$^{-1}$)} \\
    \midrule
    Off-diagonal ($H_{12}$)       & 13.52 \\
    State-energy shift ($\alpha$) & 1.30  \\
    \bottomrule
  \end{tabular}
\end{table}

%% TABLE 6 — Simulation protocol workflow
\begin{table}[htbp]
\centering
\small
\caption[Simulation protocol workflow]{\textbf{Simulation protocol workflow.} {\normalfont Comprehensive minimization, relaxation, and FEP simulation protocol. The system was gradually prepared through restrained and unrestrained minimization, followed by stepwise relaxation before production FEP calculations.}}
\label{tab:combined_protocol}
\setlength{\tabcolsep}{8pt}
\renewcommand{\arraystretch}{1.3}
\begin{tabular}{@{}lccc@{}}
\toprule
\textbf{Step/Process} & \textbf{Steps} & \textbf{Time (ps)} & \textbf{Remarks} \\
\midrule
\multicolumn{4}{@{}l}{\textit{Restrained Minimization}} \\[2pt]
\quad Steps 1--5 & 100--500  & 0.01--0.05 & Light restraints \\
\quad Step 6     & 50,000    & 25.0       & Heavy atoms (10 kcal·mol$^{-1}$·\AA$^{-2}$) \\
\addlinespace
\multicolumn{4}{@{}l}{\textit{Unrestrained Minimization}} \\[2pt]
\quad Steps 7--11 & 100--500 & 0.01--0.05 & No restraints \\
\quad Step 12     & 50,000   & 25.0       & No restraints \\
\addlinespace
\multicolumn{4}{@{}l}{\textit{Relaxation}} \\[2pt]
\quad Heat-up (1--3)       & 2,000--5,000 & 2.0--5.0 & Gradual restraint release \\
\quad Equilibration (4--8) & 50,000       & 50.0     & Progressive release \\
\quad Final                & 10,000       & 10.0     & Fully equilibrated \\
\addlinespace
\multicolumn{4}{@{}l}{\textit{Production}} \\[2pt]
\quad FEP (51 windows) & 20,000 each & 1,020 total & 1 fs timestep \\
\bottomrule
\end{tabular}
\end{table}